\title{Physics-Informed Knowledge Distillation for Real-Time State-of-Charge Estimation of $\text{LiFePO}_4$ Batteries on Automotive Microcontrollers}

\author{Hamid~Daneshvar,~\IEEEmembership{Student~Member,~IEEE,}
	and~Masoud~Masih-Tehrani,~\IEEEmembership{Member,~IEEE}%
	\thanks{Manuscript received August 20, 2026. This work was supported in part by the Vehicle, Fuel, and Environment Research Institute (VFERI) and the School of Automotive Engineering at Iran University of Science and Technology (IUST). \textit{(Corresponding author: Masoud Masih-Tehrani.)}}%
	\thanks{Hamid Daneshvar and Masoud Masih-Tehrani are with the School of Automotive Engineering, Iran University of Science and Technology (IUST), Tehran 16846-13114, Iran (e-mail: hamid\_daneshvar@auto.iust.ac.ir; masih@iust.ac.ir).}}

\markboth{IEEE Transactions on Transportation Electrification,~Vol.~XX, No.~XX,~August~2026}%
{Daneshvar and Masih-Tehrani: Physics-Informed Knowledge Distillation for Real-Time SOC Estimation}

\maketitle

\begin{abstract}
	Precise state-of-charge (SOC) estimation in Lithium Iron Phosphate ($\text{LiFePO}_4$ or LFP) batteries remains an enduring industrial challenge for automotive battery management systems (BMS) due to the ultra-flat open-circuit voltage (OCV) plateau ($\text{d}V_{\text{ocv}}/\text{d}\text{SOC} \approx 0.1\text{--}0.5\text{ mV/\%}$) and protracted non-linear solid-phase relaxation dynamics. While heavy data-driven ensembles and electrochemical partial differential equation (PDE) solvers capture these path-dependent transitions, their computational complexity and excessive memory footprints prohibit deterministic real-time execution on resource-constrained automotive microcontrollers. In this study, we propose \textit{MicroPhys-BMS}, an end-to-end Physics-Informed Machine Learning Knowledge Distillation (PIML-KD) framework that compresses complex non-linear electro-thermal relaxation dynamics into an ultra-compact student Micro-MLP architecture ($12 \to 32 \to 16 \to 1$, 961 parameters) for ISO 26262 ASIL-D compliant microcontrollers. The framework extracts a 12-dimensional domain-guided physical feature vector encoding Arrhenius reaction kinetics, dynamic ohmic overpotential, logarithmic solid-phase diffusion scaling, and differential relaxation gradients. A 300-tree gradient-boosted decision tree (LightGBM) teacher supervises the student network via a composite physics-calibrated objective combining adaptive Huber distillation, task-level mean squared error, and quadratic physical boundary constraints enforcing $\text{SOC} \in [0.0, 1.0]$. Validated against extensive Stanford LFP operational telemetry, the proposed student Micro-MLP achieves a nominal $25^\circ\text{C}$ case-split MAE of 0.454\% and RMSE of 1.236\% ($R^2 = 0.9986$), outperforming the LightGBM teacher (MAE = 0.488\%). Under unseen cell generalization tests, the model records SOC MAEs of 2.48\% on cell YX07 and 3.35\% on cell YX08, reducing estimation errors by 42.3\% and 22.1\% relative to the Stanford baseline benchmark target (4.30\%). For short-term rest horizons, the framework enables robust Coulomb counting re-initialization at 30~s (MAE = 1.170\%) and converges to 0.339\% at 300~s. Furthermore, the model demonstrates high resilience against analog front-end (AFE) voltage noise (MAE = 1.336\% under $\pm 5.0\text{ mV}$ perturbation) and multi-temperature shifts (MAE = 2.689\% at 35$^\circ$C), with 95.94\% of cumulative error distribution remaining within $\le 2.0\%$. Synthesized into a MISRA-C:2012 compliant static C header without dynamic heap allocation (\texttt{malloc}), the deployed model requires only 3.75~kB Flash, 192~Bytes SRAM, and an execution latency of $\approx 12\ \mu\text{s}$ on an ARM Cortex-M4 at 160~MHz, achieving over 99.7\% memory compression compared to standard tree ensembles.
\end{abstract}

\begin{IEEEkeywords}
	Automotive microcontrollers, battery management system (BMS), knowledge distillation (KD), lithium iron phosphate (LFP) battery, physics-informed machine learning (PIML), state-of-charge (SOC).
\end{IEEEkeywords}